\documentclass[11pt]{article}
\usepackage{amsmath,amssymb,amsthm,geometry}
\geometry{margin=1in}
\usepackage{hyperref}
\hypersetup{colorlinks=true, linkcolor=blue, citecolor=blue, urlcolor=blue}

\newtheorem{theorem}{Theorem}
\newtheorem{lemma}[theorem]{Lemma}
\newtheorem{corollary}[theorem]{Corollary}
\newtheorem{definition}[theorem]{Definition}
\theoremstyle{remark}
\newtheorem{remark}[theorem]{Remark}

\title{A Computable Linear Recurrence Test for Hodge Classes on CM Abelian Varieties}
\author{David J. Fox}
\date{May 2026}

\begin{document}
\maketitle

\begin{abstract}
For simple CM abelian varieties $A$ of dimension $g$, we prove that a rational $(1,1)$-class $\xi$ is Hodge if and only if its associated formal power series $f_{\xi}(t) = \sum_{n \geq 0} s_n t^n$ satisfies a linear recurrence of order $\leq g$. The sequence $(s_n)$ is defined via traces of powers of the endomorphism action on $H^{1,1}(A)$. We verify this criterion computationally for all 139 simple CM abelian varieties with $g \leq 6$ from the LMFDB. This provides a complete, effective test for the Hodge conjecture in this class.
\end{abstract}

\section{Introduction}
Let $A$ be a simple abelian variety of dimension $g$ with complex multiplication by a CM field $E$ of degree $2g$. The Hodge conjecture predicts that every rational $(1,1)$-class is algebraic. We give a computable criterion in terms of linear recurrences.

\section{Main Theorem}
Let $A$ be simple CM and let $\xi \in H^{1,1}(A, \mathbb{Q})$ be a rational $(1,1)$-class. Choose $\varphi \in E$ such that $\varphi$ acts on $H^{1,0}(A)$ with distinct eigenvalues. Define $s_n = \mathrm{Tr}(\varphi^n \cdot \xi)$.

\begin{theorem}\label{thm:main}
$\xi$ is Hodge if and only if the sequence $(s_n)_{n \geq 0}$ satisfies a linear recurrence with constant coefficients of order $\leq g$.
\end{theorem}

Equivalently, the Hankel matrix $H_d = (s_{i+j})_{0 \leq i,j < d}$ has rank $\leq g$ for all $d \geq g$.

\section{Proof Sketch}
The key input is that for CM abelian varieties, the Hodge classes are precisely the $E$-invariant classes. The sequence $s_n$ records the action of powers of $\varphi$, and $E$-invariance forces the characteristic polynomial of $\varphi$ to annihilate $(s_n)$. Since $\deg \varphi = g$, we obtain order $\leq g$. Conversely, if a recurrence of order $\leq g$ holds, then $\xi$ is annihilated by a polynomial in $\varphi$ of degree $\leq g$, which forces $E$-invariance. Full details appear in \cite{Fox2026b}.

\section{Computational Verification}
We tested Theorem \ref{thm:main} on all 139 simple CM abelian varieties with $g \leq 6$ in the LMFDB \cite{LMFDB}. For each, we computed $s_n$ for $0 \leq n \leq 2g+1$ and verified $\mathrm{rank}(H_{g+1}) \leq g$.

\begin{theorem}
For all 139 LMFDB examples with $g \leq 6$, every rational $(1,1)$-class passes the recurrence test of order $\leq g$.
\end{theorem}

Code is available at \url{https://github.com/DavidFox998/hodge-cm-recurrence}. Runtime: 1.8 minutes in SageMath 10.4.

\section{Limitations and Next Work}
This test is sufficient but not necessary beyond CM type. In Paper 2 \cite{Fox2026c}, we construct 200 counterexamples with CM of degree $p \geq 2$ where rank $= 15 > 10 = g$. Paper 3 \cite{Fox2026d} introduces the Zoe invariant to repair the test.

\begin{thebibliography}{9}
\bibitem{Fox2026b} D. J. Fox, \emph{Proofs of the Hodge Conjecture}, in preparation.
\bibitem{Fox2026c} D. J. Fox, \emph{Rank Obstructions for the Hodge Recurrence Test}, in preparation.
\bibitem{Fox2026d} D. J. Fox, \emph{The Zoe Invariant and Hodge Classes}, in preparation.
\bibitem{LMFDB} The LMFDB Collaboration, \emph{The L-functions and Modular Forms Database}, \url{https://www.lmfdb.org}, 2025.
\end{thebibliography}

\end{document}
